%%%%% slides template
% Giovanni Ramirez (ramirez@ecfm.usac.edu.gt)
% 20200828
% Require: images/ecfmByN.png images/usacByN.jpg
%
%
% Use this option to print notes only
% \documentclass[xcolor=dvipsnames,handout,notes=only]{beamer}%
%
% Use this option to print notes and slides
% \documentclass[xcolor=dvipsnames,handout,notes=show]{beamer}%
%
% Use this option to get a printer-friendly version
% \documentclass[xcolor=dvipsnames,handout]{beamer}%
%
% Unse this option to get slides only
\documentclass[xcolor=dvipsnames,presentation]{beamer}%
%
%%%%%


%%%%% theme and coloring
% a list of themes can be found here https://hartwork.org/beamer-theme-matrix/
% I like a simple Malmoe with some modifications and color-schemes like
% dolphin for a white background, beetle for a gray background or dove for
% plain slides
%
% theme definition
\usetheme{Malmoe}
%
% definition of the presentation mode: full slides
\mode<presentation> {
% coloring for slides other options: beetle, dove
  \usecolortheme{dolphin}
% to include a slide with the table of contents
 \AtBeginSection[] {
  \begin{frame}
    \tableofcontents[currentsection]
  \end{frame} }
}
%
% definition of the handout mode: printer-friendly
\mode<handout>{
% coloring: dove is white
  \usecolortheme{dove}
% this is to print 2 slides in 1 page, other values are allowed
  \usepackage{pgfpages}
  \renewcommand\textbullet{\ensuremath{\bullet}}
  \pgfpagesuselayout{2 on 1}[letterpaper,border shrink=5mm]
% set the fontsize of the notes
  \setbeamerfont{note page}{size=\scriptsize}
% disable the slide with the table of contents
  \AtBeginSection[]{}
}
%%%%%

%%%%% personalisation
% clear all headers
\setbeamertemplate{headline}{}
% disable navigation buttons
\setbeamertemplate{navigation symbols}{}
% configure the blocks used for equations, theorems, etc
\setbeamertemplate{blocks}[rounded][shadow=true]
% clear the footline
\setbeamertemplate{footline}{}
% set the footline
\setbeamertemplate{footline}{
  \hbox{%
% left block with the shortname, it is set in the \author command
    \begin{beamercolorbox}[wd=.40\paperwidth,ht=2.25ex,dp=1ex,center]%
      {author in head/foot}%
      \usebeamerfont{author in head/foot}\insertshortauthor
    \end{beamercolorbox}%
% center block with the short title, it is set in the \title command
    \begin{beamercolorbox}[wd=.50\paperwidth,ht=2.25ex,dp=1ex,center]%
      {title in head/foot}%
      \usebeamerfont{title in head/foot}\insertshorttitle
    \end{beamercolorbox}%
% right block with the number of slides
    \begin{beamercolorbox}[wd=.10\paperwidth,ht=2.25ex,dp=1ex,right]%
      {date in head/foot}%
      \usebeamerfont{date in head/foot}
      \insertframenumber{} / \inserttotalframenumber\hspace*{1em}
    \end{beamercolorbox}}%
  \vskip0pt%
}
%%%%%

%%%%% packages
%
% language, decimal point mark
\usepackage[spanish,es-nodecimaldot]{babel}%
%
% to use only T1 scalable fonts
\usepackage[T1]{fontenc}%
\let\Tiny=\tiny% this is a proper configuration of the tiny font size
%
% to use UTF coding, other options latin1
\usepackage[utf8]{inputenc}%
%
\usepackage{array}
% special math symbols
\usepackage{latexsym,amsfonts,amsmath}%
%
% to include graphics
\usepackage{graphicx}%
%
% to use tiks: nice to draw diagrams
\usepackage{tikz}%
\usetikzlibrary{arrows,shapes,positioning}
%
% to use SI units and physics notation
\usepackage[squaren]{SIunits}%
\usepackage{physics}%
\usepackage{multirow}
\usepackage{xcolor}
%
% to use hypenation
\usepackage{ragged2e}%
\let\raggedright=\RaggedRight%
%
\usepackage{dsfont}
% to configure the hyperref to set metadata in the PDF file
\hypersetup{%
%pdfpagemode=FullScreen,% to open the file in fullscreen
breaklinks={true},% to allow links in more than one line
%pdfstartview={Fit},% to set the slide to fit the screen
%pdfpagemode=FullScreen,% to set it to full screen
pdfauthor={Giovanni Ramirez Garcia (ramirez@ecfm.usac.edu.gt)},% author name
pdftitle={Quantum Optics aplied to Quantum Computing},% title of the talk
pdfsubject={Keynote for the talk at Lat0 group},% subject of the talk
pdfkeywords={quantum optics, quantum computing, quantum algorithms, quantum
  circuits}% keywords to categorise the contents
}
%%%%%

%%%%% Talk data
%
%title
\title[Canales cuánticos PCE]% this is the short title
{\bf Mapeos proyectivos en sistemas de qubits}% this is the long title
%
% author
\author[José Alfredo de León (ECFM-USAC)]% this is the short name
{\bf José Alfredo de León${}^1$, 
Carlos Pineda${}^2$,\\
David Dávalos${}^2$, 
Alejandro Fonseca${}^2$}% this is the long name
%
% affiliation
\institute{${}^1$Escuela de Ciencias Físicas y Matemáticas, USAC \newline%
${}^2$Instituto de Física, UNAM}%
%
% date and information about the conference
\date{I Congreso Guatemalteco de Física\\09 de julio de 2021}
%%%%%


\begin{document}
% % % % % % % % % % % % % % % % % % % % % % % % % % % % % % % % % % % % % %
% Title and Outline
% % % % % % % % % % % % % % % % % % % % % % % % % % % % % % % % % % % % % %
\begin{frame}[plain]
  \titlepage %
  \begin{tikzpicture}[x=1mm,y=1mm,overlay,remember picture]
    \pgftransformshift{\pgfpointanchor{current page}{center}}
    \node[inner sep=0pt] (usac) at (-50,0) %
    {\includegraphics[width=15mm]{images/usacByN}};
    \node[inner sep=0pt] (ecfm) at (50,0) %
    {\includegraphics[width=15mm]{images/ecfmByN}};
  \end{tikzpicture}
\end{frame}

\begin{frame}{Outline}
  \tableofcontents[hidesubsections] 
\end{frame}


% % % % % % % % % % % % % % % % % % % % % % % % % % % % % % % % % % % % % %
% Introduction and Motivation
% % % % % % % % % % % % % % % % % % % % % % % % % % % % % % % % % % % % % %
\section{Sistemas cuánticos abiertos}
\label{sec:Intro}

\begin{frame}{Sistemas cuánticos abiertos}{¿Qué son?}

	\only<1>{
	Considerar a los sistemas cuánticos como aislados de su entorno
	es una buena aproximación...
  
  \vfill

	\centering
	\includegraphics[height=6cm]{images/cerrados.pdf}
  }
  
  \only<2>{
  .. no obstante, los sistemas cuánticos reales están abiertos a interacciones
  con su entorno.
  
  \vfill

	\centering
	\includegraphics[height=6cm]{images/abiertos.pdf}
  }
\end{frame}


\begin{frame}{Sistemas cuánticos abiertos}%
  {Matriz de densidad $\rho$}\pause
  
  Si un sistema se encuentra en alguno de los estados $\ket{\psi_i}$ 
  con probabilidad $p_i$, la matriz de densidad del sistema se define como
  \begin{align*}
  \rho \equiv \sum_{i}p_i\dyad{\psi_i}{\psi_i}.
  \end{align*} \pause
  
  \vfill
  
	Una matriz $\rho$ es una matriz de densidad asociada 
	a un ensamble de estados $\{p_i,\ket{\psi_i}\}$ si y sólo si
	\begin{enumerate}
	\item $\Tr (\rho)=1$,
	\item $\rho\geq 0$.
	\end{enumerate}
	
	\vfill
\end{frame}


\begin{frame}{Sistemas cuánticos abiertos}%
  {Canales cuánticos}\pause
  
  \only<2>{
	Una operación lineal $\mathcal{E}$ que actúa sobre 
	la matriz de densidad $\rho$ 	como 
	\begin{align*}
	\mathcal{E}(\rho)=\rho'
	\end{align*}
	es un canal cuántico si
	\begin{enumerate}
	\item Preserva las características de la matriz de densidad,
	\item Es una operación completamente positiva.
	\end{enumerate}
	}
\end{frame}	

\begin{frame}{Completa positividad}{¿Para qué o qué?}
	\only<1-2>{
	\begin{center}
	\includegraphics[height=6cm]{images/CP1}
	\end{center}
	}	
	\only<3>{
	\begin{center}
	\includegraphics[height=6cm]{images/CP2}
	\end{center}	
	}
	
	\onslide<2-3>{
	No es suficiente que 	$\mathcal{E}(\rho_A)\geq0$. }\onslide<3>{También se tiene que cumplir 
	$\mathcal{E}\otimes \mathds{1}(\rho_{total})\geq0$. \vspace{.5cm}}
\end{frame}

\begin{frame}{Completa positividad}{Definición}
	\onslide<2>{
	Una operación $\mathcal{E}$ es completamente positiva 
	si y sólo si 
	\begin{align*}
	\mathcal{E}\otimes \mathds{1}(\rho_{\mathcal{E}})\geq 0,
	\end{align*}
	con $\rho_{\mathcal{E}}$ el estado maximamente entrelazado
	de un sistema compuesto por dos copias del sistema
	sobre el que actúa $\mathcal{E}$.	\vfill

	\textbf{Para recordar:} 
	$\mathcal{E}\otimes \mathds{1}(\rho_{\mathcal{E}})$ recibe 
	el nombre de la \alert{matriz de Choi}~de~$\mathcal{E}$.
	}
\end{frame}

% % % % % % % % % % % % % % % % % % % % % % % % % % % % % % % % % % % % % %
% Qubits
% % % % % % % % % % % % % % % % % % % % % % % % % % % % % % % % % % % % % %

\section{Qubits}

\begin{frame}{Qubits}
	Un qubit es un sistema cuántico de dos niveles.
	
	\begin{columns}
	\begin{column}{0.5\textwidth}
	\only<2>{
		\begin{align*}
			\rho &= \frac{1}{2}\sum _{i=0}^3 r_i\sigma_i \\
			\rho &= \frac{\mathds{1}+r_1\sigma_x+r_2\sigma_y+r_3\sigma_z}{2},
		\end{align*}	
		donde $\qty(r_1,r_2,r_3)$ se conoce como el vector de Bloch.
	}
	\end{column}
	\begin{column}{0.5\textwidth}  %%<--- here
  		\begin{center}
    		\includegraphics[width=0.7\textwidth]{images/bloch}      
     \end{center}
	\end{column}
	\end{columns}
\end{frame}

\begin{frame}{Canal cuántico de 1 qubit}{Bit-flip}	
	Geométricamente,
	\begin{center}
	\begin{tabular}{m{2.5cm} m{1cm} m{2.5cm}}
		\includegraphics[width=3cm]{images/bloch-ball}
		& \hfill $\longmapsto$
		& \includegraphics[width=3cm]{images/bit-flip}
	\end{tabular}
	\end{center}\pause
	
	El canal cuántico de bit-flip 	actúa sobre las componentes de Pauli
	de la matriz de densidad de 1 qubit como
	\begin{align*}
	\qty(1,r_1,r_2,r_3) \longmapsto \qty(1,r_1,\qty(1-p)r_2,\qty(1-p)r_3).
	\end{align*}
\end{frame}

\begin{frame}{Qubits}{$n$ qubits}
	En general, la matriz de densidad de $n$ qubits puede escribirse 
	en la base de productos tensoriales de las matrices de Pauli como
	\begin{align*}
		\rho = \frac{1}{2^n}\sum _{j_1,\ldots,j_n=0}^3 r_{j_1,\ldots,j_n}
		\sigma_{j_1}\otimes \ldots
		\otimes\sigma_{j_n},
	\end{align*}
	con $r_{0,\ldots,0}=1$ $(\Tr\rho=1)$. \pause
	Llamaremos \alert{componentes de Pauli} a las $r_{j_1,\ldots,j_n}$.
\end{frame}

% % % % % % % % % % % % % % % % % % % % % % % % % % % % % % % % % % % % % %
% Theoretical framework
% % % % % % % % % % % % % % % % % % % % % % % % % % % % % % % % % % % % % %
\section{Operaciones PCE}
\label{sec:Theory}

\begin{frame}{Motivación (1/2)}{Bit-flip con $p=0.5$} \pause
	Geométricamente, actúa como 
	\begin{center}
	\begin{tabular}{m{2.5cm} m{1cm} m{2.5cm}}
		\includegraphics[width=3cm]{images/bloch-ball}
		& \hfill $\longmapsto$
		& \includegraphics[width=3cm]{images/bitFlip-toX}
	\end{tabular}
	\end{center}
	
	\vfill
	
	En términos de las componentes de Pauli, el bit-flip para $p=0.5$
	actúa como
	\begin{align*}
	\qty(1,r_1,r_2,r_3)\longmapsto \qty(1,r_1,0,0).
	\end{align*}
\end{frame}


\begin{frame}{Motivación (2/2)}{Mapeo de la esfera de Bloch a un disco} \pause
	Consideremos la operación $\mathcal{E}$ que deforma
	la esfera de Bloch como
	\begin{center}
	\begin{tabular}{m{2.5cm} m{1cm} m{2.5cm}}
		\includegraphics[width=3cm]{images/bloch-ball}
		& \hfill $\longmapsto$
		& \includegraphics[width=3.5cm]{images/diskXY}
	\end{tabular}
	\end{center}
	\begin{align*}
	\qty(1,r_1,r_2,r_3)\longmapsto\qty(1,r_1,r_2,0).
	\end{align*}
	\pause
		
	$\mathcal{E}\otimes \mathds{1}(\rho_{\mathcal{E}})$
	no es una matriz positiva (i.e. $\mathcal{E}$ no es completamente positiva), 
	entonces $\mathcal{E}$ no es un canal cuántico	. \vspace{.5cm}
\end{frame}


\begin{frame}{Operaciones PCE}{Definición y planteamiento del problema}\pause
	Una operación PCE (\textit{Pauli component erasing}) 
	es una operación lineal que actúa sobre las
	componentes de Pauli de una matriz de densidad $\rho$ 
	de $n$ qubits como
	\begin{align*}
	r_{j_1,\ldots,j_n}\longrightarrow \tau_{j_1,\ldots,j_n}r_{j_1,\ldots,j_n},
	\end{align*}
	con $\tau_{j_1,\ldots,j_n}=0,1$; excepto $\tau_{0,\ldots,0}=1$.
	Es decir, borra o deja invariantes las componentes de Pauli. \pause
	
	\hfill
	\includegraphics[height=1.5cm]{images/1qubit-pce-example001}
	\hspace{1.8cm}
	\includegraphics[height=1.5cm]{images/2qubit-pce-example001}
	\hfill
	\includegraphics[height=2cm]{images/3qubit-pce-example001}
	\hspace{1cm} \pause

	\vfill	
	
	\colorbox{Yellow}{\textbf{¿Cómo es la caracterización del subconjunto}}
 	\colorbox{Yellow}{\textbf{de las operaciones PCE que son canales cuánticos?}}
% 	\colorbox{Yellow}{\textbf{que	son canales cuánticos?}} 
\end{frame}


\begin{frame}{Canales cuánticos PCE}{Diagonalización de la matriz de Choi} \pause
Los eigenvalores de la matriz de Choi de una operación PCE 
de $n$ qubits son
\begin{align*}
%\qty(\bigotimes_{}^na) \vec{\tau} \geq \vec{0},\\
\vec{\lambda} = 
\underbrace{\qty(a\otimes \ldots\otimes a)}_{n\text{ veces}} \vec{\tau}
\end{align*}
con
\begin{align*}
a=\mqty(1&1&1&1\\1&1&-1&-1\\1&-1&1&-1\\1&1&-1&-1)
\end{align*}
y $\vec{\tau}$ un vector de $4^n$ componentes.
\end{frame}


\begin{frame}{Canales cuánticos PCE}{$2^k$, pero también importa el orden}
Los canales cuánticos PCE dejan $2^k$ componentes
invariantes. Sin embargo, no sólo importa cuántas 
componentes de Pauli, sino también cuáles. 

\vfill

\textbf{2 qubits:} \pause
\begin{columns}
\begin{column}{0.5\textwidth}
No son canales cuánticos:
\vspace{.3cm}
\begin{center}
\includegraphics[height=1.8cm]{images/2qubit-pce-BadoneHalf}
\end{center}
\vspace{1.6cm}
\end{column} \pause
\begin{column}{0.5\textwidth}
Todos los canales cuánticos PCE que borran 15 componentes:
\begin{center}
\includegraphics[height=3.8cm]{images/2qubit-pce-oneHalf}
\end{center}
\end{column}
\end{columns}
\end{frame}

\begin{frame}{Canales cuánticos PCE}{¿Cuántos?}
Los canales cuánticos parecen seguir una regla `espejo' para 
el número de canales PCE según la cantidad de componentes 
de Pauli invariantes. \pause

\vfill

\textbf{3 qubits:}
\begin{table}[]
\centering
%\begin{tabular}{c|lllll}
%\multirow{3}{*}{\textbf{canales PCE}} &                                &                                & \multicolumn{1}{c}{35}         &                                &                                 \\
%                                                      &                                & \multicolumn{1}{c}{15}         &                                & \multicolumn{1}{c}{15}         &                                 \\
%                                                      & \multicolumn{1}{c}{1}          &                                &                                &                                & \multicolumn{1}{c}{1}           \\
%\hline 
%\textbf{componentes}                                  & \multicolumn{1}{c}{\textbf{1}} & \multicolumn{1}{c}{\textbf{2}} & \multicolumn{1}{c}{\textbf{4}} & \multicolumn{1}{c}{\textbf{8}} & \multicolumn{1}{c}{\textbf{16}}
%\end{tabular}
%\end{table}
%
%\begin{table}[]
%\begin{tabular}{c|lllll}
%\multirow{3}{*}{\textbf{operaciones PCE}} &                                &                                & \multicolumn{1}{c}{455}        &                                &                                 \\
%                                                      &                                & \multicolumn{1}{c}{15}         &                                & \multicolumn{1}{c}{6435}       &                                 \\
%                                                      & \multicolumn{1}{c}{1}          &                                &                                &                                & \multicolumn{1}{c}{1}           \\
%\hline
%\textbf{componentes}                                  & \multicolumn{1}{c}{\textbf{1}} & \multicolumn{1}{c}{\textbf{2}} & \multicolumn{1}{c}{\textbf{4}} & \multicolumn{1}{c}{\textbf{8}} & \multicolumn{1}{c}{\textbf{16}}
%\end{tabular}
\begin{table}[]
\footnotesize
\begin{tabular}{cccccccc}
\multirow{4}{*}{\textbf{canales PCE}} &            &            &            & 1395       &             &             &             \\
                                                      &            &            & 651        &            & 651         &             &             \\
                                                      &            & 63         &            &            &             & 63          &             \\
                                                      & 1          &            &            &            &             &             & 1           \\
\textbf{componentes}                                  & \textbf{1} & \textbf{2} & \textbf{4} & \textbf{8} & \textbf{16} & \textbf{32} & \textbf{64}
\end{tabular}
\end{table}
\end{table} \pause

\begin{table}[]
\footnotesize
\begin{tabular}{cccccccc}
\multirow{4}{*}{\textbf{operaciones PCE}} &            &            &            & $5.5\times 10^7$ &                     &                         &             \\
                                          &            &            & 39711      &             & $1.2\times 10^{14}$ &                         &             \\
                                          &            & 63         &            &             &                     & $9.2\times 10^{17}$ &             \\
                                          & 1          &            &            &             &                     &                         & 1           \\
\textbf{componentes}                      & \textbf{1} & \textbf{2} & \textbf{4} & \textbf{8}  & \textbf{16}         & \textbf{32}             & \textbf{64}
\end{tabular}
\end{table}

\end{frame}

\onslide<1->{
\begin{frame}{Canales cuánticos PCE}{Generadores}
Un canal cuántico PCE puede descomponerse como concatenación 
de los elementos de un subconjunto generador.
}

\onslide<2->{
\vfill
\textbf{Generadores canales PCE de 2 qubits:}

\begin{center}
\includegraphics[height=4cm]{images/2qubits-generators}
\end{center}
}

\end{frame}




\begin{frame}{Ejemplo: 2 qubits}
\onslide<1->{
Generadores:
\begin{center}
\includegraphics[width=8cm]{images/2qubits-generators}
\end{center}
}

\onslide<2->{
\hfill
\begin{tabular}{m{0.8cm} m{.1cm} m{0.8cm} m{.05cm} m{0.8cm}}
\includegraphics[height=1.1cm]{images/E1}  
& $\bigcirc$ 
& \includegraphics[height=1.1cm]{images/E4} 
& $=$
& \includegraphics[height=1.1cm]{images/E1E4} 
\end{tabular} \hspace{1.3cm}
}

\onslide<3->{
\hfill
\begin{tabular}{m{0.8cm} m{.1cm} m{0.8cm} m{.1cm} m{0.8cm} m{.1cm} m{0.8cm}}
\includegraphics[height=1.1cm]{images/E3}
& $\bigcirc$ 
& \includegraphics[height=1.1cm]{images/E1}  
& $\bigcirc$ 
& \includegraphics[height=1.1cm]{images/E4} 
& $=$
& \includegraphics[height=1.1cm]{images/E3E1E4} 
\end{tabular} \hspace{1.3cm}
}

\onslide<4->{
\hfill
\begin{tabular}{m{0.8cm} m{.1cm} m{0.8cm} m{.1cm} m{0.8cm} m{.1cm} m{0.8cm} m{.1cm} m{0.8cm}}
\includegraphics[height=1.1cm]{images/E11}
& $\bigcirc$ 
& \includegraphics[height=1.1cm]{images/E3}
& $\bigcirc$ 
& \includegraphics[height=1.1cm]{images/E1}  
& $\bigcirc$ 
& \includegraphics[height=1.1cm]{images/E4} 
& $=$
& \includegraphics[height=1.1cm]{images/E11E3E1E4} 
\end{tabular} \hspace{1.3cm}
}
\end{frame}

\begin{frame}{Hipótesis}
Sea $PCE_n$ al conjunto de canales cuánticos PCE de $n$
qubits. Existe un conjunto $\Gamma_n\subset PCE_n$ cuyos
elementos son necesarios y suficientes para generar al resto
de los elementos en $PCE_n$. \vfill \pause

Sea $\mathcal{E}_j\in \Gamma_n$ y $\Phi \in PCE_n$. Un canal 
PCE de $n$ qubits puede descomponerse como
\begin{align*}
\Phi = \underbrace{\mathcal{E}_{j_1}\circ \mathcal{E}_{j_2}\circ \ldots
\circ \mathcal{E}_{j_l}}_{\text{máximo }2^n}.
\end{align*}\vfill \pause

Los elementos de $\Gamma_n$ son los canales PCE
que proyectan la matriz de densidad $\rho$, escrita en la base de Pauli,
a un subespacio que tiene la mitad de la dimensión del espacio sobre 
el cual actúa $\Phi$.
\end{frame}

% % % % % % % % % % % % % % % % % % % % % % % % % % % % % % % % % % % % % %
% Let's have good manners
% % % % % % % % % % % % % % % % % % % % % % % % % % % % % % % % % % % % % %
\begin{frame}
  \begin{center}
    ¡Muchas gracias!

    \vfill

    Contacto:

    José Alfredo de León

    deleongarrido.jose@gmail.com
  \end{center}
\end{frame}

\end{document}

%%% Local Variables:
%%% mode: latex
%%% TeX-master: t
%%% End:
